\chapter{Conclusion}

\section{Overview of the Study}
\textit{This study provides a multi-dimensional comparative evaluation of Monolithic and Microservice architectures, 
transitioning from theoretical debates to evidence-based insights within the specific operational rigors of a Learning Management System.} 

This research conducted a rigorous and empirically grounded comparative analysis of two dominant 
software architectural paradigms — Monolithic and Microservice — within the specific operational 
context of a Learning Management System (LMS). Rather than relying on theoretical conjecture or 
domain-agnostic benchmarks, the research was designed to capture the behavioral characteristics 
of each architecture under workload conditions representative of real-world LMS deployments: concurrent 
student enrollment, course browsing during peak semester loads, and write-intensive submission workflows. 

By constructing two functionally identical implementations of the same LMS business logic — one as a 
single-process NestJS application and the other as a service containerized ecosystem coordinated via 
Nginx and RabbitMQ — and subjecting both to an identical suite of stress tests using \texttt{k6}, the 
study produced a multi-dimensional dataset spanning throughput, response time, CPU utilization, 
resource efficiency (RUE), fault tolerance, and development-phase availability.


\section{Key Findings and Insights}
\textit{The synthesis of experimental data reveals a complex interplay between architectural topology and system performance, where the advantages of service decomposition are balanced against the overhead of distributed coordination.}

The experimental results yield several interconnected findings that define the performance-cost-resilience trade-off landscape under LMS-specific conditions.

\subsection{Throughput and Response Time: The Concurrency Crossover}
\textit{The capacity of each architecture to handle concurrent traffic is not static but is governed by a critical crossover threshold that determines the efficiency of request processing.}



The throughput analysis identifies a context-dependent performance hierarchy governed by a crossover threshold at approximately \textbf{1,000 Virtual Users}. 
\begin{itemize}
    \item \textbf{Sub-threshold (Low Load):} The Monolithic architecture maintains a measurable advantage due to the absence of inter-service network hops. At single-user load, the Monolith's response time (12.8 ms) is over 2.6 times faster than the Microservice baseline (34.1 ms).
    \item \textbf{Post-threshold (High Load):} The hierarchy inverts as the Monolith reaches event-loop saturation. At 2,000 VUs, Microservice achieves 579.17 RPS compared to the Monolith’s 430.94 RPS — a \textbf{34.4\% superiority}. Under instantaneous burst patterns (Scenario 2), this advantage expands to \textbf{63\%} at 1,500 VUs, demonstrating the parallel absorption capacity of distributed containers.
\end{itemize}

However, the Microservice write-path exhibited a critical vulnerability: a P99 response time of 30,002 ms at moderate loads (200--500 VUs), indicating transaction failures due to inter-service coordination timeouts. This highlights a reliability risk for high-stakes enrollment tasks in distributed architectures.

\subsection{CPU Utilization and Resource Efficiency: The Infrastructure Tax}
\textit{The economic viability of an architecture is determined by its ability to convert computational power into successful user transactions, a ratio that fluctuates significantly between idle and peak states.}

The analysis confirms the existence of an \textbf{"Infrastructure Tax"} in Microservice. Components like the RabbitMQ broker and Nginx gateway establish a static resource floor, with RabbitMQ consuming 40--60\% CPU even at idle. 

Despite this, the \textbf{Resource Utilization Efficiency (RUE)} analysis reveals:
\begin{itemize}
    \item \textbf{Normal Load (200 VUs):} Microservice is significantly more efficient, with an $RUE_{CPU}$ of 1,148.09 compared to the Monolith’s 573.22. It processed 4.65x more requests per unit of CPU.
    \item \textbf{Peak Load (1,000 VUs):} The ranking reverses. The Monolith’s $RUE_{CPU}$ is 4.58x higher than Microservice, as the latter struggles with the exponential overhead of cross-service communication and container orchestration.
\end{itemize}

\subsection{Fault Tolerance and Development-Phase Resilience}
\textit{Architectural resilience is evaluated not only by runtime stability but also by the system's ability to isolate failures and maintain availability throughout the continuous development lifecycle.}

Empirical examination of 148 commit events reveals that Microservice achieved a system availability of \textbf{43.31\%} versus \textbf{35.14\%} for the Monolith during development. This is attributed to the "Bounded Blast Radius" property; whereas a configuration error in the Monolith's single \texttt{Dockerfile} paralyzed the entire system, errors in Microservice (e.g., Nginx pathing) resulted only in localized service outages, allowing the rest of the ecosystem to remain operational.

\section{Contextual Superiority}
\textit{Selecting the appropriate architecture requires an objective alignment between the institution's operational requirements, expected user scale, and technical infrastructure maturity.}

A central finding of this research is that neither architecture is universally superior. The optimal choice depends on operational scale and organizational maturity:

\begin{itemize}
    \item \textbf{Small to Medium Scale (<500 Concurrent Users):} The \textbf{Monolithic} architecture is the rational choice. It offers superior write-path reliability, lower maintenance complexity (lower MTTR), and higher resource efficiency during traffic spikes.
    \item \textbf{Large Scale (>2,000 Concurrent Users):} \textbf{Microservice} are decisively superior. Their ability to maintain high throughput and isolate faults during peak academic events outweighs their higher baseline infrastructure cost.
    \item \textbf{The Crossover Range (500--2,000 Users):} This "ambiguous zone" requires selection based on DevOps maturity. Organizations should consider the \textbf{Monolithic} as an intermediate strategy to retain operational simplicity while preparing for future decomposition.
\end{itemize}

\section{Final Remarks}
\textit{This study concludes that architectural design is an ongoing process of managing trade-offs, where data-driven insights provide the necessary clarity to balance performance, cost, and resilience.}

This research contributes a rigorous empirical framework that challenges both the naive adoption of Microservice and the conservative retention of Monoliths. The findings demonstrate that while Microservice provide a "Stability Dividend" and better fault isolation, they impose a heavy "Coordination Tax" during peak loads. 

Ultimately, architectural selection for an LMS should be a data-driven decision based on the specific "Performance Knee" and "Efficiency Crossover" points identified in this study. In system design, every decision is a calculated trade-off; this work provides the empirical instruments to make those trade-offs with precision and clarity.